% !TeX encoding = UTF-8
\documentclass[variant=apuntes,cover=institutional,mode=screen]{ua-engineering-v6-article}
% !TeX encoding = UTF-8
\UASetUniversity{Universidad Austral}
\UASetFaculty{Facultad de Ingeniería}
\UASetDepartment{Departamento de Computación}
\UASetCampus{Pilar}
\UASetCareer{Ingeniería Informática}
\UASetCommission{Comisión A}
\UASetProfessor{Equipo docente}
\UASetSemester{Segundo cuatrimestre 2026}
\UASetCourse{Sistemas Distribuidos}
\UASetDate{2026}


\UASetClassNumber{16}
\UASetClassTitle{Defensa final ante comité}
\UASetDocKind{Apuntes}

\title{Clase 16 --- Defensa final de sistema ante comité: presentación, objeciones, revisión crítica y rúbrica de evaluación}
\author{\UAProfessor}
\date{\UADate}

\begin{document}
\maketitle

\section{Introducción}

El cierre natural de un curso serio de sistemas distribuidos no consiste simplemente en recordar mecanismos, sino en demostrar la capacidad de integrarlos en una propuesta coherente y defenderla frente a objeciones técnicas.

Esta clase toma todo lo visto y lo reordena desde la perspectiva de una \textbf{defensa final ante comité}. Eso implica cambiar nuevamente el nivel de análisis:

\begin{itemize}
\item ya no se trata solo de saber qué es un algoritmo o patrón;
\item se trata de justificar por qué una arquitectura concreta es razonable;
\item se trata de responder críticamente bajo preguntas adversariales;
\item se trata de reconocer límites, no de ocultarlos.
\end{itemize}

\begin{quote}
Una buena defensa arquitectónica no intenta parecer perfecta. Intenta ser coherente, explícita y técnicamente honesta.
\end{quote}

\section{Qué significa defender un sistema}

Defender un sistema no es describir su diagrama ni enumerar herramientas. Tampoco es un ejercicio de venta.

Una defensa sólida debe responder, de manera integrada, estas preguntas:

\begin{itemize}
\item ¿qué problema resuelve el sistema?;
\item ¿qué restricciones del dominio importan?;
\item ¿qué fallas se asumen como relevantes?;
\item ¿qué garantías se consideran indispensables?;
\item ¿qué arquitectura se propone y por qué?;
\item ¿qué costo se acepta al tomar esa decisión?;
\item ¿cómo se opera, observa y mejora el sistema?;
\item ¿qué límites y riesgos permanecen?
\end{itemize}

\section{La estructura mínima de una defensa}

Una secuencia razonable para una defensa final es la siguiente.

\subsection{1. Problema y contexto}

Debe quedar claro:
\begin{itemize}
\item quién usa el sistema;
\item qué flujo de negocio resuelve;
\item qué operaciones son críticas;
\item qué escala, latencia o distribución geográfica se esperan.
\end{itemize}

Una arquitectura sin problema explícito tiende a ser decorativa.

\subsection{2. Requisitos y restricciones}

Deben distinguirse:
\begin{itemize}
\item requisitos funcionales;
\item requisitos no funcionales;
\item restricciones organizacionales o operativas;
\item restricciones regulatorias o de negocio.
\end{itemize}

\subsection{3. Modelo de falla}

Una defensa madura no parte del supuesto de red perfecta. Debe decir qué fallas considera:

\begin{itemize}
\item latencia variable;
\item timeouts;
\item pérdida;
\item duplicación;
\item partición;
\item crash de nodos;
\item degradación de dependencias;
\item errores humanos y operacionales.
\end{itemize}

\subsection{4. Garantías}

Aquí deben explicitarse las garantías buscadas:
\begin{itemize}
\item consistencia fuerte o eventual;
\item durabilidad;
\item semántica de entrega;
\item auditabilidad;
\item resiliencia;
\item SLOs.
\end{itemize}

\subsection{5. Arquitectura propuesta}

No basta con dibujar cajas. Deben justificarse:
\begin{itemize}
\item fronteras del sistema;
\item bounded contexts;
\item ownership de datos;
\item modelo de comunicación;
\item uso de RPC y/o eventos;
\item mecanismos de coordinación o compensación;
\item estrategia de almacenamiento y operación.
\end{itemize}

\subsection{6. Trade-offs aceptados}

Toda propuesta seria debe decir qué costo paga:
\begin{itemize}
\item más latencia;
\item más complejidad operativa;
\item menor disponibilidad bajo ciertas fallas;
\item más reconciliación posterior;
\item mayor costo de infraestructura.
\end{itemize}

\subsection{7. Límites y trabajo futuro}

Una defensa fuerte reconoce:
\begin{itemize}
\item qué parte sigue siendo frágil;
\item qué hipótesis todavía deberían validarse;
\item qué decisiones dependen de pruebas adicionales;
\item qué parte del diseño está conscientemente postergada.
\end{itemize}

\section{El marco del curso como columna vertebral}

El curso propuso reiteradamente el marco:

\[
D = (P, F, G, S, T)
\]

Esta es probablemente la forma más compacta y poderosa de estructurar una defensa final.

\subsection{Problema ($P$)}
Define el dominio y evita arquitecturas sin anclaje.

\subsection{Fallas ($F$)}
Evita el diseño idealizado y obliga a pensar en sistemas reales.

\subsection{Garantías ($G$)}
Obliga a distinguir qué propiedades importan realmente.

\subsection{Solución ($S$)}
Permite presentar la arquitectura como respuesta al problema, no como moda.

\subsection{Trade-offs ($T$)}
Introduce honestidad técnica y evita promesas incompatibles.

\section{Qué mira un comité serio}

Un comité técnico no suele premiar la propuesta con más componentes o más buzzwords. Valora especialmente:

\begin{itemize}
\item claridad conceptual;
\item consistencia interna;
\item calidad de las prioridades;
\item explicitación de supuestos;
\item capacidad de responder objeciones;
\item honestidad sobre límites.
\end{itemize}

En cambio, penaliza:
\begin{itemize}
\item slogans;
\item respuestas totalizantes;
\item promesas incompatibles;
\item desconocimiento del modelo de falla;
\item ausencia de operación o de observabilidad;
\item improvisación evasiva.
\end{itemize}

\section{Objeciones típicas}

En una defensa final, algunas objeciones aparecen casi inevitablemente.

\subsection{“¿Por qué no un monolito?”}

Esta pregunta obliga a justificar que la complejidad distribuida resuelve un problema real y no una preferencia ideológica.

Una buena respuesta debe hablar de:
\begin{itemize}
\item límites del dominio;
\item autonomía de equipos;
\item escalado diferencial;
\item ownership de datos;
\item costo aceptado al distribuir.
\end{itemize}

\subsection{“¿Por qué no consistencia fuerte en todo?”}

Esta objeción obliga a mostrar comprensión del costo de:
\begin{itemize}
\item coordinación;
\item latencia;
\item disponibilidad bajo partición;
\item complejidad operacional.
\end{itemize}

\subsection{“¿Qué pasa bajo partición?”}

Es una pregunta central porque revela si la propuesta existe en el mundo real o solo en el camino feliz.

\subsection{“¿Cómo operás esto?”}

Sin respuesta a observabilidad, SLOs, incidentes, ownership y postmortems, la arquitectura queda incompleta.

\section{Cómo responder bien}

Una buena respuesta ante objeciones técnicas tiene una forma reconocible:

\begin{itemize}
\item reconoce el supuesto relevante;
\item explica qué se priorizó;
\item explicita el costo aceptado;
\item muestra por qué esa elección tiene sentido para el dominio.
\end{itemize}

Lo que no debe hacerse:
\begin{itemize}
\item negar el trade-off;
\item responder con slogans;
\item inventar garantías que el sistema no tiene;
\item ocultar complejidad operacional;
\item prometer simultáneamente propiedades incompatibles.
\end{itemize}

\section{Revisión crítica de la propia propuesta}

Una defensa madura puede y debe incluir autocrítica explícita:
\begin{itemize}
\item dónde están los puntos frágiles;
\item qué parte es más costosa;
\item qué parte dependería de validación empírica;
\item qué parte podría rediseñarse si cambiara el dominio.
\end{itemize}

Esto no debilita la propuesta. La fortalece, porque demuestra criterio técnico y ausencia de ingenuidad.

\section{Rúbrica de evaluación}

Una rúbrica razonable para evaluar una defensa final puede incluir:

\subsection{1. Formulación del problema}
\begin{itemize}
\item claridad del contexto;
\item relevancia del dominio;
\item precisión de requisitos.
\end{itemize}

\subsection{2. Modelo de falla}
\begin{itemize}
\item explicitación de fallas;
\item realismo de supuestos;
\item tratamiento de casos adversos.
\end{itemize}

\subsection{3. Calidad del diseño}
\begin{itemize}
\item coherencia de la arquitectura;
\item adecuación del modelo de comunicación y datos;
\item consistencia entre objetivos y mecanismos.
\end{itemize}

\subsection{4. Trade-offs}
\begin{itemize}
\item explicitación de costos;
\item madurez en la priorización;
\item honestidad sobre límites.
\end{itemize}

\subsection{5. Operación}
\begin{itemize}
\item observabilidad;
\item resiliencia;
\item confiabilidad;
\item incidentes y aprendizaje.
\end{itemize}

\subsection{6. Defensa oral}
\begin{itemize}
\item claridad expositiva;
\item calidad de respuestas;
\item manejo de objeciones;
\item consistencia argumentativa.
\end{itemize}

\section{Qué distingue una defensa sobresaliente}

Una defensa sobresaliente no es la que promete más. Es la que logra:
\begin{itemize}
\item formular un problema con precisión;
\item proponer una arquitectura coherente;
\item alinear mecanismos con garantías;
\item hacer visibles sus límites;
\item responder preguntas hostiles sin perder rigor.
\end{itemize}

\section{Conclusión}

El cierre del curso no está en memorizar patrones, sino en alcanzar la capacidad de integrar, argumentar y defender una arquitectura distribuida completa. Ese es, probablemente, el signo más claro de madurez conceptual en el área.

\begin{uakeybox}
Saber sistemas distribuidos significa poder convertir complejidad en una propuesta coherente y defendible bajo preguntas difíciles.
\end{uakeybox}


\section{Síntesis razonada de la clase}
Esta unidad debe leerse como parte de una cadena de decisiones de diseño y no como un catálogo de mecanismos. Para estudiar el tema con profundidad conviene reconstruir cuatro preguntas: \emph{qué problema aparece}, \emph{qué garantía se desea}, \emph{qué mecanismo la aproxima} y \emph{qué costo introduce}. En cada ejemplo debe identificarse además el modelo de falla implícito.

\begin{uakeybox}
Una respuesta técnicamente madura no termina al nombrar un algoritmo o patrón: declara sus supuestos, la propiedad que preserva y el escenario en el que deja de ser suficiente.
\end{uakeybox}

\section{Bibliografía guiada y ruta de profundización}
Para profundizar esta clase se recomienda: Relectura transversal de Kleppmann, papers del curso y material SRE.

La lectura debe hacerse con preguntas concretas: (1) ¿qué modelo de sistema asume el autor?, (2) ¿qué propiedad demuestra o persigue?, (3) ¿qué falla o anomalía motiva el mecanismo?, y (4) ¿qué trade-off queda abierto?

\section{Conexión con el Trabajo Final Integrador}
El grupo debe señalar qué parte de su sistema utiliza los conceptos de esta unidad, justificar por qué son necesarios y documentar una alternativa descartada. Cuando el contenido todavía no corresponda a la etapa vigente, debe registrarse como hipótesis o decisión pendiente.

\section{Autoevaluación conceptual}
\begin{enumerate}
\item ¿Cuál es el problema distribuido concreto que motiva esta clase?
\item ¿Qué propiedad de seguridad y qué propiedad de progreso están en juego?
\item ¿Qué supuesto de red, tiempo o falla resulta decisivo?
\item ¿Qué trade-off aceptaría en un sistema real y cuál no aceptaría en un dominio crítico?
\item ¿Cómo observaría experimentalmente que la solución se comporta como espera?
\end{enumerate}

\section{Profundización autocontenida v9}
\subsection{Problema, límite e invariante}
El núcleo de esta clase es presentación ante comité, objeciones, revisión crítica, rúbrica y defensa oral individual. Para estudiarlo de manera autónoma conviene separar tres planos. Primero, el \emph{problema observable}: qué comportamiento incorrecto puede ver un cliente o un operador. Segundo, el \emph{límite del modelo}: qué información, sincronía o coordinación no está disponible. Tercero, el \emph{invariante}: qué propiedad no debe violarse aunque el sistema pierda progreso temporalmente.

El modelo de fallas relevante incluye presentación descriptiva, respuesta evasiva, límite no reconocido y rúbrica inconsistente. Una solución debe describirse como protocolo y no solo como nombre de tecnología: actores, estados, mensajes, condición de éxito, condición de aborto y estado visible después de una interrupción. La evidencia esperada es presentación de 12-15 minutos, prototipo, preguntas y evaluación con rúbrica.

\subsection{Método de análisis}
Ante un caso nuevo, siga esta secuencia:
\begin{enumerate}
\item Identifique actores, dato o recurso compartido y frontera de confianza.
\item Escriba una ejecución nominal y una ejecución adversarial.
\item Distinga seguridad (lo malo no ocurre) de progreso (algo bueno finalmente ocurre).
\item Declare qué supuesto permite avanzar: mayoría, deadline, sincronía parcial, operación idempotente, almacenamiento durable u otro.
\item Compare una alternativa más simple y una más fuerte; registre qué métrica o experimento haría cambiar la decisión.
\end{enumerate}

\subsection{Caso transferible}
Considere una plataforma de reservas que recibe operaciones duplicadas, pierde conectividad entre regiones y debe sostener un camino crítico sin ocultar el estado real al usuario. Aplique el mecanismo de esta clase a una única decisión concreta. Una respuesta madura no intenta resolver todo: delimita la operación, formula la garantía y explica la degradación esperada cuando el supuesto central deja de cumplirse.

\subsection{Errores de razonamiento frecuentes}
\begin{itemize}
\item confundir una propiedad con el producto que suele implementarla;
\item describir el camino feliz sin recorrer una falla intermedia;
\item afirmar ``exactly-once'', ``alta disponibilidad'' o ``consistencia'' sin definir alcance observable;
\item medir solo throughput aunque la hipótesis trate sobre semántica o recuperación;
\item presentar la complejidad añadida como beneficio sin compararla con la alternativa más simple.
\end{itemize}

\section{Lectura guiada}
Rúbrica del curso, ADRs del proyecto y bibliografía técnica utilizada por cada equipo.

Durante la lectura, registre: modelo de sistema, propiedad perseguida, contraejemplo que motiva el mecanismo, costo principal y una condición bajo la cual no lo elegiría. Estas cinco anotaciones convierten la bibliografía en una herramienta de diseño y no en una lista para memorizar.

\section{Aplicación al Trabajo Final Integrador}
El equipo debe producir un ADR breve con la decisión asociada a esta clase. Debe contener: contexto, dos alternativas, supuesto decisivo, garantía elegida, trade-off, evidencia pendiente y fecha de revisión. La evidencia mínima esperada para esta unidad es: presentación de 12-15 minutos, prototipo, preguntas y evaluación con rúbrica.

\section{Autoevaluación avanzada}
\begin{enumerate}
\item ¿Puedo explicar la clase sin nombrar primero una tecnología?
\item ¿Puedo construir una ejecución que viole la garantía si el mecanismo se configura mal?
\item ¿Puedo distinguir seguridad, progreso y experiencia del usuario?
\item ¿Puedo proponer una medición que valide la propiedad, no solo el rendimiento?
\item ¿Puedo defender cuándo una solución más simple sería suficiente?
\end{enumerate}

\end{document}
